\subsection{run.var}
\texttt{run.var} holds the Meso-NH simulation parameters proper (timestep, nesting, microphysics,
domain geometry), the plotting height/limit parameters, and the configuration of the DIAG/FICVAL
extraction (which parameters, at which points/domains/heights, are exported as ASCII for plotting and
for the short-timescale hand-off). Most variable names mirror the corresponding Meso-NH namelist
keyword. This file must be sourced after \texttt{times.var}.

\subsubsection{Simulation parameters}
\begin{lstlisting}
NUM_PROC=integer   : number of MPI processes for the simulation. [current: 64]
NAM_SIMU="string"  : 5-character CEXP (Meso-NH experiment name). [current: "${DAY}_${MONTH}"]
NAM_SEG="string"   : 5-character CSEG (Meso-NH segment name). [current: "y${YEAR}"]
NMODEL_VAR=integer : number of models/domains, normally equal to NESTING.
XTSTEP_VAR=seconds : outermost-domain Meso-NH timestep. [current: 12]
NDTRATIO_ARRAY[n]  : timestep ratio of nested level n+1 with respect to its parent
                     (element [0]=1 always). [current: 1,2,3,4]
NDAD_ARRAY[n]      : parent-level index of nested level n+1 (element [0]=1 always).
                     [current: 1,1,2,3]
XRIMKMAX[n]        : maximum lateral-relaxation coefficient (1/s) for nested level
                     n+1; kept hand-tuned rather than auto-computed (a commented-out
                     auto-generation formula is kept in the file for reference, but is
                     not used -- it was found to lead to instability).
                     [current: 0.0006666, 0.0020000, 0.0013333]
CCLOUD_VAR="string": microphysical scheme. [current: "KESS"]
CPRESOPT_VALUE="string": pressure solver. [current: "ZRESI"]
NUMOUTTIMES=integer: number of Meso-NH (.lfi) output times. [current: 15]
OUTTIMESTEP=seconds: interval between outputs; defaults to EXITSEC from times.var --
                     changing one without the other causes inconsistencies.
\end{lstlisting}
\texttt{XSEGLEN} (total simulated seconds) and the \texttt{XFMOUT[]} array (the list of output times
themselves) are computed automatically from \texttt{NUMOUTTIMES} and \texttt{OUTTIMESTEP}; they should
not be edited directly.

\subsubsection{Spatial configuration and coordinates}
\begin{lstlisting}
XSNG1="string", YSNG1="string": grid coordinates (I,J) of the first extraction point
                     (p1.dat). [current: "61","61"]
XSNG2="string", YSNG2="string": grid coordinates (I,J) of the second extraction point
                     (p2.dat). [current: "60","61"]
SEESTPOINT="string" : starting point index for seeing integration (see_p1_v.dat).
                     [current: "2"]
DOM_FIRST="string"  : nesting level of the first extraction point (1=outer). [current: "3"]
DOM_SECOND="string" : nesting level of the second extraction point. [current: "3"]
SUPCOMP_VAR="F"     : enables supplementary computations; always "F" in operation.
MESODIAG_VAR="F"    : "F"=normal run, "T"=DIAG-only mode (internal use).
NKMAX_VALUE="string": number of vertical levels. [current: "54"]
ZDZGRD_VALUE="string": near-ground vertical mesh size (m). [current: "20"]
ZDZTOP_VALUE="string": near-top vertical mesh size (m). [current: "600"]
ZZMAX_STRGRD_VALUE="string": altitude (m) separating the two constant-stretching layers.
                     [current: "3500"]
ZSTRGRD_VALUE="string": constant stretching (%) in the lower layer. [current: "20"]
ZSTRTOP_VALUE="string": constant stretching (%) in the upper layer. [current: "00"]
\end{lstlisting}

\subsubsection{Plot parameters}
\begin{lstlisting}
PLOTBASEH1, PLOTBASEH2 : lower plot height limits (m). [current: 20, 600]
PLOTTOPH1, PLOTTOPH2   : upper plot height limits (m). [current: 600, 20000]
CN2H2=meters           : intermediate height limit for CN2/wind/RH plots. [current: 2000]
MRH2=meters            : intermediate height limit for mixing-ratio plots. [current: 5000]
CN2_WRESCALE_H=meters  : height above which the wind-gradient correction is applied to
                          CN2 profiles. [current: 700]
CN2_ALPHA, CN2_BETA    : coefficients of that correction. [current: 1., 1.5]
DOM_FIRST_DX,
DOM_SECOND_DX=km       : grid spacing used for 2D maps (must match the PGD grid).
                          [current: 0.5, 0.5]
WINDSPEEDHIGHLIM=m/s   : threshold on the average near-ground (K=4) wind speed, above
                          which backup.sh still uploads the wind figures/maps even when
                          full-figure upload is otherwise disabled -- see
                          Section~\ref{sec:usage}. [current: 8.5]
\end{lstlisting}
Changes to this block generally require matching manual edits to \texttt{plot\_python.sh}.

\subsubsection{DIAG/FICVAL extraction parameters}
These control which parameters are exported by the DIAG/FICVAL phase (\texttt{postproc.sh}), and are
also the parameters whose file layout \texttt{plot\_python.sh} depends on -- changing them may require
matching edits there.
\begin{lstlisting}
EXTRACTLIST_DOM_FIRST,
EXTRACTLIST_DOM_SECOND="string": space-separated list of diaprog keywords for 2D ground
                     maps / point extraction, for DOM_FIRST and DOM_SECOND respectively.
                     [current: DOM_FIRST carries the full list (wind, SNG, SCI, ISO, CHT,
                      ECH and 200/10227/9276/6677 m wind levels), DOM_SECOND empty]
EXTRACTLIST_MULTIHEIGHT="string": keywords extracted as a function of height.
                     [current: "SNG_STAR_SUP CHT_STAR_SUP"]
PVLIST_DOM_FIRST,
PVLIST_DOM_SECOND="string": diaprog keywords for vertical profile extraction.
                     [current: both empty]
STARTPV_I[0],STARTPV_J[0]: grid point for the DOM_FIRST-related vertical profile.
                     [current: 61,61]
STARTPV_I[1],STARTPV_J[1]: grid point for the DOM_SECOND-related vertical profile.
                     [current: 51,51]
CVLIST_DOM_FIRST,
CVLIST_DOM_SECOND="string": diaprog keywords for vertical cuts. [current: both empty]
STARTCV_I/J[0], ENDCV_I/J[0]: start/end grid points of the DOM_FIRST-related vertical
                     cut. [current: (2,61) to (121,61)]
STARTCV_I/J[1], ENDCV_I/J[1]: start/end grid points of the DOM_SECOND-related vertical
                     cut. [current: (2,51) to (101,51)]
NUMLEVELDIAG=integer: number of nesting levels on which DIAG runs. [current: 1]
DIAGLEVELARRAY[n]   : which nesting level DIAG runs on, for each of the NUMLEVELDIAG
                     entries. [current: [0]=$DOM_FIRST, [1]=$DOM_SECOND (only [0] used
                     since NUMLEVELDIAG=1)]
DIAGINFHEIGHTS[n],
DIAGSUPHEIGHTS[n]   : space-separated lists of lower/upper integration heights for
                     DIAGLEVELARRAY[n]; every INF-SUP combination is computed.
                     [current: [0]="$PLOTBASEH1 $PLOTBASEH2" / "$PLOTTOPH1 $PLOTTOPH2"]
DOMP1=1, DOMP2=2    : fixed identifiers for the two extraction-point domains; used
                     elsewhere in plot.sh's domain selection logic and should not be
                     changed.
\end{lstlisting}
